../cictikz/src/cictikz/data/tex/cictikz_preamble.tex

% Small signal model including the bulk transconductance g_s*v_sb (the
% back-gate). Note the arrow direction: the g_s source points up, as in
% the original hand drawing.

\begin{circuitikz}[american, thick, transform shape]

  ../cictikz/src/cictikz/data/tex/cictikz_lib.tex

  \def\ytop{2.6}
  \def\xgm{2.2}
  \def\xgs{4.4}
  \def\xrds{5.6}
  \def\xterm{6.8}

  % Gate terminal, floating
  \draw (0.09,\ytop) circle (0.09);
  \node[anchor=east] at (-0.1,\ytop) {$v_g$};
  \node[anchor=north] at (0.1,{\ytop-0.25}) {$+$};
  \node[anchor=west, color=blue] at (-0.15,1.3) {$v_{gs}$};
  \node[anchor=south] at (0.1,0.25) {$-$};

  % Transconductance and back-gate transconductance
  \draw (\xgm,\ytop) to [cI, l_=$g_m v_{gs}$] (\xgm,0);
  \draw (\xgs,0) to [cI, l=$g_s v_{sb}$] (\xgs,\ytop);

  % Output resistance
  \draw (\xrds,0) -- (\xrds,0.5) \vresistor{$r_{ds}$} -- (\xrds,2.6);

  % Rails to the terminals
  \draw (\xgm,\ytop) -- (\xterm,\ytop);
  \fill (\xgs,\ytop) circle (0.075);
  \fill (\xrds,\ytop) circle (0.075);
  \draw (\xterm+0.09,\ytop) circle (0.09);
  \node[anchor=west] at ({\xterm+0.3},\ytop) {$v_d$};

  \draw (0.0,0) -- (\xterm,0);
  \fill (\xgm,0) circle (0.075);
  \fill (\xgs,0) circle (0.075);
  \fill (\xrds,0) circle (0.075);
  \draw (\xterm+0.09,0) circle (0.09);
  \node[anchor=west] at ({\xterm+0.3},0) {$v_s$};

\end{circuitikz}
\end{document}
